\cvsection{Research Experience}
\cvrole{Ph.D. Thesis: Efficient Video Action Detection}{2025 – 2026}{CISE Department, University of Florida}{Gainesville, FL}
\begin{itemize}
\item Designed a causal, online (streaming) temporal action-detection model on state-space (Mamba) backbones for higher temporal accuracy and lower latency in long videos.
\item On THUMOS, the model reaches 73.3 streaming mAP with a ResNet-50 backbone and 76.2 with DINOv2, matching the strongest verified strictly-causal detector at roughly 60\% of its parameters and 2.5–2.9× its throughput, and runs in real time at 37.7 FPS.
\item Built a unified framework to benchmark online, sliding-window, and multimodal video models on THUMOS and ActivityNet across mAP, latency, throughput, memory, and FLOPs, with distributed multi-GPU training on a SLURM-managed HPC cluster.
\item First-author manuscript in preparation.
\end{itemize}
\cvrole{ECOLE Program – CReLeRI Team (DARPA-funded)}{July 2023 – May 2025}{Ruiz HCI Lab and DSR Lab, CISE Department, University of Florida}{Gainesville, FL}
\begin{itemize}
\item Developed and integrated UI and image/video analysis APIs and trained/tested action segmentation models for an explainable long-video analysis system; first author of its ACM Multimedia 2025 demo paper.
\item Built the system's long-video analysis web tool: an HTML/JS/CSS front end and a FastAPI backend with Celery and Redis background tasks, containerized with Docker, with GPU compute offloaded to UF's HiPerGator cluster.
\item Applied vision-language and LLM-based models, using LLaMA to parse Qwen-generated captions in the final system.
\end{itemize}
\cvrole{PTG Program – ENKIx Team (DARPA-funded)}{July 2023 – December 2024}{Ruiz HCI Lab, CISE Department, University of Florida}{Gainesville, FL}
\begin{itemize}
\item Trained and tested action recognition models, including fine-tuning the video-language model LaViLa, for a real-time AR task guidance system; achieved 34\% F1-score in task recognition (mid-range results among DARPA performers) and conducted usability and cognitive load evaluations.
\end{itemize}
\cvrole{Animal Behavior Recognition (NSF-funded)}{August 2020 – July 2023}{Graphics, Imaging, and Light Measurement Lab, CISE Department, UF}{Gainesville, FL}
\begin{itemize}
\item Investigated adversarial approaches for modeling 3D motion and behavior in videos for applications in biomedical research.
\item Conducted a survey of animal action recognition algorithms, trained and tested existing SOTA methods, and began implementing novel semi-supervised adversarial algorithms.
\end{itemize}
\cvrole{Rhinoplasty Research Group}{May 2019 – May 2020}{Computer Science Department, Florida Polytechnic University}{Lakeland, FL}
\begin{itemize}
\item Conducted a survey of over fifty research papers, documenting the facial landmarks and measurements relevant to the field of rhinoplasty.
\item Collaborated with other students to develop a web interface for importing, exporting, and displaying feature point coordinates and measurement values onto 3D models of the human face.
\item Achieved 95\% agreement with Blender software and outperformed 2D methods in inter-/intra-rater reliability and usability metrics.
\end{itemize}
